Face à l'émergence de crises de plus en plus complexes et systémiques : climatiques, sécuritaires, technologiques ou informationnelles, les \AssociationsUnionIHEDN{} peuvent apporter une contribution utile, à condition d'en préciser le cadre et les limites. Leur valeur ajoutée se situe principalement en amont et en aval des crises, à travers des actions de préparation, de sensibilisation, d'analyse et de retour d'expérience.

Cependant, l'examen de l'annuaire 2022 montre que les associations ne regroupent qu'une part minoritaire de la communauté des auditeurs de l'IHEDN et que les données disponibles renseignent peu sur l'activité réelle, les compétences ou les disponibilités, malgré le potentiel de compétences et de connaissances de l'ensemble de ces auditeurs.

Par ailleurs, l'utilisation effective de ces compétences suppose préalablement une connaissance fine et objectivée des profils, des expertises et des disponibilités des membres. Dans cette perspective, la mise en place d'un questionnaire visant à cartographier les compétences au sein de ces associations apparaît comme un objectif pertinent et justifié.

Une telle démarche doit néanmoins s'inscrire dans un cadre méthodologique rigoureux et s'appuyer sur un plan prévisionnel précis, afin de garantir la fiabilité des données recueillies, la pertinence des analyses produites et, \textit{in fine}, la crédibilité du dispositif proposé.

Le présent rapport expose ce raisonnement et propose une méthode de travail prenant en compte à la fois les ressources disponibles et les limites propres aux \AssociationsUnionIHEDN.